\documentclass{beamer}

\usetheme{Szeged}
\usecolortheme{beaver}

\usepackage[utf8]{inputenc}
\usepackage[spanish, mexico]{babel}
\usepackage{amsmath}
\usepackage{amssymb}
\usepackage{txfonts}
\usepackage{subcaption}
\usepackage[T1]{fontenc}

\usepackage[absolute,overlay]{textpos}
\usepackage{physics,bm}
\usepackage{enumerate}
\usepackage{rsfso}
\usepackage{graphicx,url,twoopt,natbib}

\setbeamertemplate{blocks}[rounded][shadow]
\setbeamercolor{block body}{bg=example text.fg!10}
\setbeamercolor{block title}{bg=example text.fg!20}
\setbeamercolor{block body alerted}{bg=alerted text.fg!10}
\setbeamercolor{block title alerted}{bg=alerted text.fg!20}
\setbeamertemplate{caption}[numbered]

\title[Visualización de orbitales del átomo de Helio mediante Hartree-Fock]{VISUALIZACIÓN DE ORBITALES DEL ÁTOMO DE HELIO POR EL MÉTODO DE HARTREE-FOCK}
\author[J. Jara]{Jovanny Jara}
\institute [Optimización y Programación GPU]{
Facultad de Ciencias \\
Licenciatura en Física \\
Universidad de Valparaíso
}
\date{\today}

\begin{document}

\begin{frame}
    \titlepage  
\end{frame}

\begin{frame}{Objetivo del proyecto}
    El objetivo principal del proyecto es resolver la ecuación de Schrödinger para el átomo de Helio para obtener sus energías permitidas y visualizar dichos estados.
\end{frame}

\begin{frame}{Problema a resolver}
    Atomo de Helio: dos electrones + núcleo $\to$ Problema de n cuerpos. \ \alert{No hay solución exacta!} 
    \pause
    \\\vspace{4mm} Necesitamos método numérico. $\to$ \alert{Hartree-Fock}
\end{frame}

\begin{frame}{Método de Hartree-Fock}
    El Hamiltoniano del Helio (en unidades atómicas):
    \begin{equation}
        \hat{\mathrm{H}} = \frac{1}{2}\nabla_1^2 + \frac{1}{2}\nabla_2^2 + \frac{Z}{\text{r}_1} + \frac{Z}{\text{r}_2} + \frac{1}{\text{r}_{12}}
    \end{equation}
    \pause 
    \\ El último término es la repulsión entre los electrones. \alert{Nos acopla el problema.}
\end{frame}

\begin{frame}{Método de Hartree-Fock}
    Idea central del método:
    \begin{itemize}
        \item Cada electrón se mueve en un potencial promedio (campo medio) generado por el otro electrón
    \end{itemize}
    \alert{Consecuencia:} Podemos aproximar la función de onda total como el producto de dos orbitales de un solo cuerpo.
    \begin{equation}
        \boxed{\Psi(\text{r}_1, \text{r}_2) \approx \phi(\text{r}_1)\phi(\text{r}_2)}
    \end{equation}
\end{frame}

\begin{frame}{Método de Hartree-Fock}
    Esto nos reduce el problema a un problema de un cuerpo efectivo, donde pasamos ahora a resolver la ecuacion de Fock:
    \begin{equation}
        \hat{F}\phi(\text{r}) = \epsilon\phi(\text{r})
    \end{equation}
    Donde $\epsilon$ es la energía del orbital $\phi$ y $\hat{F}$ es el operador de Fock, que se escribe como
    \begin{equation}
        \hat{F} = \frac{1}{2}\nabla^2 + \frac{Z}{\text{r}} + V_H(\text{r}) 
    \end{equation}
    Siendo $V_H$ el potencial de Hartree, que corresponde al potencial efectivo producido por la densidad de carga del otro electron.
\end{frame}

\begin{frame}{Método de Hartree-Fock}
    El problema es que estas ecuaciones nuevas no son lineales: Dependemos de $\phi(\text{r}_2)$ para encontrar $\phi(\text{r}_1)$ y viceversa. Esta dependencia reside en el potencial de Hartree. \\\vspace{4mm}
    Si interpretamos la densidad de carga (en unidades atomicas) del otro electron como 
    \begin{equation}
        \rho = |\phi(\text{r}_j)|^2
    \end{equation}
    Entonces, el potencial de Hartree queda descrito
    \begin{equation}
        V_H(\text{r}_i) = \int \frac{\phi^*(\text{r}_j)\phi(\text{r}_j)}{\text{r}_{ij}}d\text{r}_j
    \end{equation}
\end{frame}

\begin{frame}{Método de Hartree-Fock}
    Lo que hacemos para salir del dilema:
    \begin{enumerate}[1.-]
        \item Proponemos $\phi(\text{r}_j)$
        \item Calculamos $V_H(\text{r}_i)$
        \item Resolvemos la ecuacion de Fock y obtenemos $\phi(\text{r}_i)$
        \item Calculamos $V_H(\text{r}_j)$
        \item Resolvemos la ecuacion de Fock y obtenemos un nuevo $\phi(\text{r}_j)$
        \item Repetimos desde el paso 2 hasta que $\phi(\text{r}_i)$ y $\phi(\text{r}_j)$ no cambien. \alert{Convergencia.}
    \end{enumerate}
    Esto es, en escencia, el metodo de Hartree-Fock.\\\vspace{4mm}
\end{frame}

\begin{frame}{Implementación del método}
    Ahora nos quedan dos tareas: Como implementamos el calculo de $V_H$ y la ecuacion de Fock. \\\vspace{4mm}
    Cálculo de $V_H$: Calcular integrales numéricamente es costoso, pero hay otra forma de hacerlo más simple: \alert{resolver la ecuación de Laplace.}\\
    Por que? $\to$ Podemos usar FFTs (cuFFT)
    \begin{equation}
        \nabla^2 V_H = -4\pi\rho \xrightarrow{FFT} \frac{1}{k^2}V_H = -4\pi\rho
    \end{equation}
\end{frame}

\begin{frame}{Implementación del método}
    Ecuación de Fock: Normalmente resolver un problema de valores propios requiere diagonalizar matrices (complicado) $\to$ Otro efoque: Imaginary Time Propagation (ITP).\\\vspace{4mm}
    La ec. de Schrödinger dependiente del tiempo
    \begin{equation}
        -i\hbar\pdv{\Psi}{t} = \hat{H}\Psi
    \end{equation}
    Sustituyendo $t = -i\tau$
    \begin{equation}
        \pdv{\Psi}{\tau} = -\hat{H}\Psi \implies \Psi(\tau) = e^{-\hat{H}\tau}\psi(0)
    \end{equation}
    Esta sustitución nos convierte la solucion ondulatoria a un decaimiento $\implies$ \alert{Convergencia al estado de menor energía}
\end{frame}

\begin{frame}{Implementación del método}
    El problema: El potencial (depende de r) y la energía cinética (depende de p) no conmutan.\\\vspace{4mm}

    La solución: Usar la aproximacion \alert{Split-Operator}
    \begin{enumerate}[1.-]
        \item Inicializamos una malla en 3D (con la función de prueba del HF)
        \item Multiplicamos por $e^{-V(r)\Delta\tau/2}$ (Split-Step)
        \item Multiplicamos por $e^{-k^2\Delta\tau/2}$ (aprovechamos FFT por la energía cinética)
        \item Volvemos a multiplicar por $e^{-V(r)\Delta\tau/2}$
    \end{enumerate}
    Y con esto concluye el paso de resolver la ecuación de fock en el ciclo SCF.
\end{frame}

\end{document}